\begin{codebox}
\codeline{\textcolor{cmtgray}{\#\ ==============================}}
\codeline{\textcolor{cmtgray}{\#\ 非单调推理\ (Non-monotonic\ Reasoning)}}
\codeline{\textcolor{cmtgray}{\#\ ==============================}}
\codeline{\textcolor{cmtgray}{\#\ 默认规则：设备一般是可用的，除非明确知道其发生故障}}
\codeblank{}
\codeline{Equipment(x)\ :\ available(x)\ /\ available(x)}
\codeline{\textcolor{cmtgray}{\#\ 规则解读：如果x是设备，且"x可用"这一假设是一致的}}
\codeline{\textcolor{cmtgray}{\#\ （即知识库中没有x故障的信息），则默认假设x可用。}}
\codeblank{}
\codeline{\textcolor{cmtgray}{\#\ 当知识库中添加\ hasFault(Lathe\_001)\ 后：}}
\codeline{\textcolor{cmtgray}{\#\ \(\rightarrow\)\ "Lathe\_001可用"\ 的假设不再成立}}
\codeline{\textcolor{cmtgray}{\#\ \(\rightarrow\)\ 之前的推论被推翻（撤销推理）}}
\codeline{\textcolor{cmtgray}{\#\ \(\rightarrow\)\ 对于没有故障记录的设备，仍然默认其可用}}
\end{codebox}
